\subsection*{Source position 0}
Why are museums only open during normal business hours ? Example : URL\_0 * * I work during normal business hours . When am I supposed to go ? * * Explain like I'm five.

\paragraph{Original; $p=5.91e-14$.} Museums often stay closed during standard work hours because the staff inside need time to clean exhibits, repair damaged items, and organize new displays without getting interrupted. Imagine if a teacher had to teach a class while someone tried to put away their textbooks or fix the broken whiteboard right in front of them; it would be chaotic and dangerous for the delicate art. By closing early, workers can safely do their special cleaning jobs so visitors can see everything clearly later. While staff might work shifts, keeping the doors shut allows them to maintain the space without rushing through tasks or risking damage to priceless treasures while trying to help customers.


\paragraph{Prune 10 / direct; quality passes; $p=9.21e-09$.} To ensure the safety of delicate artifacts and the high quality of visitor experiences, museums frequently remain closed during standard operating hours to allow dedicated staff time for essential maintenance tasks such as exhibit cleaning, artifact repair, and display organization without interruption. This early closure mimics an ideal scenario where an educator could move freely to manage a classroom rather than face chaos and potential damage while simultaneously attempting to tidy up textbooks or repair a whiteboard. By keeping the doors locked, employees can perform these specialized duties safely and methodically, avoiding the rush that often leads to accidents or harm to priceless treasures, ensuring that when the facility reopens, visitors can enjoy a clear and well-maintained environment despite the staff being on rotation shifts.

\paragraph{Prune 10 / smooth; quality passes; $p=1.91e-06$.} To ensure safety and accessibility, museums frequently remain inaccessible during conventional operating times, as interior personnel require dedicated time to clean deteriorated artifacts, restore broken pieces, and arrange fresh exhibitions without compromising their efforts to customers. Consider the scenario of an educator attempting to instruct a group while a colleague simultaneously attempts to relocate students' coursebooks or mend a malfunctioning whiteboard situated directly in their path; such disruption would prove both disorderly and perilous for fragile artistic objects. By closing operations beforehand, employees can execute their specialized maintenance duties in a secure environment, guaranteeing that all patrons encounter everything in pristine condition upon the next opening. Although staff may work multiple shifts, maintaining the facility in a closed state enables them to manage daily tasks efficiently without the need to expedite procedures or endanger valuable historical treasures while simultaneously assisting visitors.

\paragraph{Prune 10 / ratio; quality passes; $p=1.17e-07$.} To ensure the safety of delicate art and the integrity of the collection, museums frequently remain closed during standard working hours, allowing staff to dedicate their full attention to essential tasks such as curating new displays, repairing damaged pieces, and cleaning exhibits without interruption. This early closure mimics the necessity of a teacher remaining in the room to arrange textbooks or repair a whiteboard while simultaneously managing a class, a situation that would otherwise be chaotic and hazardous for the fragile objects on display. By shutting the doors temporarily, workers can safely complete these specialized maintenance jobs, ensuring that when the doors reopen later, visitors encounter everything in pristine condition; although employees may be distributed across shifts, the practice of keeping the facility closed prevents the rush of duties that could compromise the preservation of priceless treasures or cause harm to customers.

\paragraph{Prune 20 / direct; quality passes; $p=3.06e-10$.} To ensure the safety of delicate artifacts and the high quality of visitor experiences, museums frequently remain closed during standard operating hours to allow dedicated staff time for essential maintenance tasks such as exhibit cleaning, artifact restoration, and the arrangement of new displays in a disturbance-free environment. This precaution is comparable to the scenario where a teacher attempts to instruct a class while simultaneously dealing with disorganized textbooks or a broken whiteboard, a situation that would create chaos and threaten the fragile items; similarly, closing early enables workers to perform their specialized cleaning duties without interruption so that artworks are presented clearly to the public upon reopening. Although employees may operate on shift schedules, maintaining facility closures prevents the rush that often occurs when the public is present, thereby eliminating the risk of accidental damage to priceless treasures or safety hazards while staff attend to customer needs.

\paragraph{Prune 20 / smooth; quality passes; $p=2.03e-02$.} To ensure safety and clarity for visitors, museums frequently remain closed during typical business hours so that employees can dedicate the necessary time to clean displays, mend broken pieces, and arrange fresh exhibits without interruption. Consider the scenario where an instructor attempts to deliver a lesson while simultaneously managing an orderly removal of books or repairing a damaged whiteboard directly in their path; such actions would create extreme chaos and pose a lethal threat to fragile artworks. By shutting the doors early, staff members can safely complete these specialized maintenance duties, guaranteeing that the collections are pristine for the public to view shortly thereafter. Although personnel may operate on rotating schedules, keeping the facility closed allows them to preserve the environment and avoid the dangers of rushing to assist customers, thereby preventing potential harm to the invaluable treasures stored within.

\paragraph{Prune 20 / ratio; quality passes; $p=2.33e-11$.} To ensure the safety of delicate art and the integrity of the collection, museums frequently remain closed during standard working hours, allowing staff to dedicate their full attention to essential tasks such as curating new displays, repairing damaged pieces, and cleaning exhibits without interruption. This early closure mimics the necessity of a teacher withdrawing from a classroom to attend to urgent repairs or organize materials, as leaving these duties to be handled alongside public access would create a chaotic and perilous environment for both visitors and the artwork. By shutting the doors, employees can perform specialized maintenance and restocking services efficiently, thereby guaranteeing that visitors are not disturbed during the process and ensuring that no priceless treasures are compromised while workers rush to accommodate requests.

\paragraph{Independent reference; quality passes; $p=3.27e-01$.} Museums frequently close during scheduled operating periods due to the necessity of thorough maintenance, including the cleaning of artifacts, the restoration of faulty objects, and the reassembly of fresh exhibit displays without the risk of accidental interruptions. Should an educator present a lecture while a student attempts to store their books or fix a malfunctioning whiteboard situated right before the class, the result would be a chaotic and hazardous experience that endangers the fragile art. By early closing times, employees can safely execute their specialized cleaning duties, thereby ensuring that all exhibits remain accessible to the public for viewing the full scope of their collections. While staff members may alternate shifts during the workday, maintaining the interior spaces by locking the doors prevents hurried task execution that could inadvertently cause damage to the precious treasures or disrupt the educational flow for customers.

\subsection*{Source position 8}
The Unibomber Someone at a Halloween party was dressed as him an I 'm not sure what he did . ( In detail ) Explain like I'm five.

\paragraph{Original; $p=2.16e-19$.} The Unibomber was a man named Ted Kaczynski who made bombs and killed many people just to be angry at everyone. A group called the Bomb Squad has to use their special superpowers to find all the places where he hid his dangerous time travel machines. Once they found them, the police arrested him so no one else would ever have to worry about getting hurt by an evil balloon animal again forever because now the world is safe from his bad dreams.


\paragraph{Prune 10 / direct; quality passes; $p=2.75e-03$.} Ted Kaczynski, known as the Unibomber, was the individual responsible for constructing explosive devices and causing numerous fatalities driven by a desire to rage against all humanity. It fell to the Bomb Squad, a specialized unit equipped with extraordinary capabilities, to locate the hazardous time travel apparatuses that Kaczynski had concealed in various locations. After successfully identifying these sites, law enforcement apprehended him to ensure that future generations would never face the threat of injury from this malevolent figure, thereby securing eternal safety for the world from his sinister nightmares.

\paragraph{Prune 10 / smooth; quality passes; $p=1.71e-03$.} Ted Kaczynski, known as the Unibomber, constructed bombs and assassinated numerous individuals purely to instill rage upon society. The Bomb Squad, a specialized unit, employs unique extraordinary abilities to locate every site where he concealed his hazardous time travel devices. Upon discovery, the police apprehended him to ensure that future generations would never face the threat of injury from a malevolent balloon animal, thereby securing the world's safety from his malign nightmares permanently.

\paragraph{Prune 10 / ratio; quality fails; $p=2.01e-04$.} Ted Kaczynski, known as the Unibomber, was an individual who constructed explosives and caused numerous fatalities solely out of malevolence toward humanity. A specialized unit referred to as the Bomb Squad employed unique capabilities to locate all sites where Kaczynski concealed his volatile inventions. Upon their discovery, law enforcement arrested him to ensure that future generations would never face the threat of being injured by such a dangerous figure, thereby guaranteeing the world's safety from his persistent negative visions forever.

\paragraph{Prune 20 / direct; quality passes; $p=1.39e-07$.} Ted Kaczynski, known as the Unibomber, constructed explosives with the sole purpose of inflicting harm and spreading anger; consequently, a specialized unit referred to as the Bomb Squad employed their unique extraordinary abilities to locate every site where he concealed these hazardous, time-traveling devices. After successfully identifying these locations, law enforcement apprehended him to ensure that future generations would never face the threat of injury from a malevolent balloon construct, thereby guaranteeing lasting safety and liberating the globe from the shadow of his malevolent fantasies.

\paragraph{Prune 20 / smooth; quality passes; $p=9.14e-04$.} Ted Kaczynski, known as the Unibomber, created explosive devices out of anger toward humanity. Consequently, the Bomb Squad possesses the ability to locate every instance where Kaczynski concealed his dangerous time travel apparatus. Upon discovery, authorities detained him to ensure that future generations would remain unharmed from a future where he continues to manifest as an evil balloon animal, thereby guaranteeing the world's safety from his persistent nightmares.

\paragraph{Prune 20 / ratio; quality fails; $p=2.60e-05$.} Ted Kaczynski, known as the Unibomber, was an individual who constructed explosives and caused numerous fatalities solely out of malevolence toward humanity. A specialized unit referred to as the Bomb Squad employed unique capabilities to locate all sites where Kaczynski concealed his volatile inventions. Upon their discovery, law enforcement arrested him to ensure that future generations would remain secure against the threat of his creations, thereby ending the era of fear associated with his mischievous balloon-like devices and the nightmares they inspired.

\paragraph{Independent reference; quality passes; $p=2.88e-01$.} Ted Kaczynski, the individual often referred to as The Unibomber, was a notorious figure known for creating weapons that led to numerous fatalities solely to express intense rage at his adversaries. A clandestine unit known as the Bomb Squad utilizes their unique superpowers to locate every instance where Kaczynski concealed his time-traveling device. Upon locating these locations, law enforcement eventually apprehended him, thereby ensuring that no other individual could ever fear the potential consequences of an evil balloon creature, with the current global security situation safeguarded by his absence from nightmares.

\subsection*{Source position 16}
Why do hospitals and psych - wards get abandoned ? Why are these places not immediately re - purposed ? And even when they are left why do people explore these places and try to get murdered ? Please explain like I'm five.

\paragraph{Original; $p=3.98e-21$.} Hospitals close when no one needs their expensive beds or special doctors anymore, like you throw away an empty toy you never play with. Re-purposing is hard because those buildings often have dangerous old parts, high floors, or laws that stop people from changing them quickly. Explorers entering these spots are actually very dangerous; they often ignore strict rules meant to keep them safe. Some sadly try to commit suicide or get hurt while lost there, but brave rules inside usually stop bad things before they can happen to anyone foolish enough to wander in alone.


\paragraph{Prune 10 / direct; quality fails; $p=5.53e-08$.} Once the demand for costly medical facilities and specialized care disappears, hospitals often stand abandoned and unused, much like a discarded toy that has never been played with. The process of repurposing these structures is frequently impeded by hazardous aging components, elevated floors, and regulatory restrictions that hinder rapid modification. While exploring such locations is extremely perilous, as individuals often disregard the strict safety protocols designed to protect them, the authorities typically maintain rigorous internal rules that frequently intervene to prevent tragedies, including fatal accidents and suicides, from befalling reckless individuals who dare to enter these neglected sites without permission.

\paragraph{Prune 10 / smooth; quality passes; $p=3.43e-03$.} Once the demand for costly medical facilities and specialist care has vanished, these hospitals often cease operations, much like discarding an unused toy. However, converting such structures proves difficult due to hazardous relics, elevated floors, and regulations that impede rapid renovation. Adventurers venturing into these locations face significant risks, frequently disregarding strict safety protocols designed to protect them. Tragically, some individuals attempt suicide or suffer injuries while abandoned there, yet inherent safety measures generally prevent these tragedies before they reach anyone foolish enough to venture in independently.

\paragraph{Prune 10 / ratio; quality passes; $p=1.37e-04$.} Hospitals frequently shut down once their costly amenities, such as expensive beds and specialized medical staff, are no longer required, much like discarding an unused toy; however, repurposing these structures presents significant obstacles because they often contain hazardous remnants, high elevations, or regulations that impede rapid renovation efforts. Consequently, adventurers exploring these locations face severe peril, frequently disregarding established safety protocols designed to protect them. Tragically, some individuals attempt to commit suicide or suffer injuries while lost in these environments, though typically, strict internal rules are in place to prevent harm from befalling anyone who enters without proper supervision.

\paragraph{Prune 20 / direct; quality fails; $p=8.48e-11$.} Once the demand for costly medical facilities and specialized care disappears, hospitals often stand abandoned and unused, much like a discarded toy that has lost all purpose. While re-using these structures presents significant obstacles due to hazardous aging components, elevated levels, and restrictive regulations that delay renovations, venturing into them remains perilous; despite the existence of rigorous safety protocols intended to prevent accidents, many intruders disregard these warnings, sometimes resulting in tragic incidents such as suicides or injuries from getting lost. Fortunately, the severe rules strictly enforced within these sites often avert such disasters before they can occur to any reckless individual foolish enough to wander there without proper caution.

\paragraph{Prune 20 / smooth; quality passes; $p=1.24e-04$.} Once again, hospitals shut their doors when they realize that the patients no longer require their costly beds or the specialized medical staff they employed, much like discarding an empty toy that has long been in disuse. This re-purposing process proves difficult because the structures frequently harbor hazardous, outdated components, feature elevated floor levels, or are governed by regulations that impede rapid renovation efforts. The prospect of exploration within these locations is perilous, as individuals often disregard stringent safety protocols designed to protect them. Tragically, some have attempted to end their lives or sustain harm while disoriented there, though the vigilant rules within these facilities frequently intervene to prevent such tragedies from befalling anyone foolish enough to venture blindly into the unknown.

\paragraph{Prune 20 / ratio; quality fails; $p=3.99e-06$.} Hospitals frequently shut down once their costly amenities, such as expensive beds and specialized medical staff, are no longer required, much like discarding an unused, empty toy. Adapting these structures is challenging due to the presence of hazardous elements, high ceilings, or regulations that impede rapid renovation. Moreover, adventurers venturing into these sites face significant peril, as they frequently disregard safety protocols designed to protect them. Tragically, some individuals attempt to self-harm or suffer injuries while lost in these locations, yet courageous policies within the facility usually intervene to prevent such harm before it befalls anyone foolish enough to wander there alone.

\paragraph{Independent reference; quality passes; $p=2.73e-01$.} Hospitals shut down when the need for expensive beds or specialized medical personnel vanishes, much like discarding an unused toy that no longer serves a purpose. Restoring these facilities requires overcoming formidable obstacles that include hazardous architectural relics, elevated structures, or regulatory frameworks restricting their alteration. Visitors to these inaccessible areas frequently encounter perilous conditions, often deviating from safety protocols that aim to prevent harm. Despite the tragic reality that some individuals risk suicide or injury while stranded in such spaces, robust internal regulations typically act as a safeguard, stopping dangerous behavior long before it escalates for those foolish enough to venture alone.

\subsection*{Source position 24}
How to pay with cash when car shopping?

\paragraph{Original; $p=8.34e-10$.} To pay with cash when car shopping, first ensure you have a clear check or verified funds available well before visiting the dealership. Bring your government-issued identification and proof of funds, such as a recent bank statement or cashier's check request, to facilitate the transaction. Negotiate your price separately from your financing paperwork if you plan to borrow, then present the exact cash amount to the salesperson at the desk. Ensure they sign a receipt for your payment and verify that the vehicle title is properly transferred to your name immediately upon completion to avoid future ownership issues.


\paragraph{Prune 10 / direct; quality passes; $p=8.16e-05$.} Before visiting a car dealership to pay in cash, secure verified funds or a confirmed check well in advance, bringing along government-issued ID alongside evidence of your money, such as a current bank statement or a cashier's check request. If you intend to finance the purchase, handle price negotiations separately from the financing documents, then provide the precise cash sum to the salesperson, who should sign a receipt confirming the transaction. It is crucial to immediately verify that the vehicle title is legally transferred to your name upon closing to prevent any future ownership complications.

\paragraph{Prune 10 / smooth; quality passes; $p=5.81e-03$.} Before visiting a car dealership to pay with cash, verify that you have sufficient funds ready or a reliable check to present, accompanied by your government-issued ID and supporting documentation like a bank statement or a cashier's check. Discuss and agree upon the specific purchase price before initiating any financing discussions, then hand over the exact cash sum at the counter to the sales associate. Once the transaction is complete, the salesperson must sign a receipt confirming the payment, and you must confirm that the vehicle's title has been officially transferred into your name on the spot to prevent any potential ownership complications.

\paragraph{Prune 10 / ratio; quality passes; $p=1.25e-05$.} Before visiting a car dealership to settle a transaction using physical currency, you must first confirm that liquid funds are accessible, preferably through a clear check or verified account balance, and arrive at the location with your government-issued ID alongside documentation proving those funds, such as a current bank statement or a cashier's check request. Should you intend to finance the purchase, discuss the final price independently of the loan documents before meeting the salesperson, then hand over the precise cash sum at the counter. Immediately after the deal closes, ensure the salesperson signs a receipt confirming the payment and that the vehicle's title is correctly transferred into your name to prevent any subsequent ownership complications.

\paragraph{Prune 20 / direct; quality passes; $p=8.16e-05$.} Before visiting a car dealership to pay in cash, secure verified funds or a confirmed check well in advance, bringing along government-issued ID alongside evidence of your money, such as a current bank statement or a cashier's check request. If you intend to finance the purchase, handle price negotiations separately from the financing documents, then provide the precise cash sum to the salesperson, who should sign a receipt confirming the transaction. It is crucial to immediately verify that the vehicle title is legally transferred to your name upon closing to prevent any future ownership complications.

\paragraph{Prune 20 / smooth; quality passes; $p=3.60e-03$.} Before visiting a car dealership to pay with cash, verify that you have sufficient funds ready or a reliable check to present, accompanied by your government-issued ID and supporting documentation like a bank statement or a cashier's check. Discuss and agree upon the specific purchase price before any financing documents are signed, and then hand the exact cash figure to the sales staff, ensuring they provide you with a signed receipt. Finally, immediately confirm that the vehicle's title has been officially transferred into your name to prevent any future complications regarding ownership.

\paragraph{Prune 20 / ratio; quality passes; $p=1.31e-03$.} Before visiting a car dealership to settle a transaction using physical currency, you must first confirm the availability of sufficient cash or a verified check in your possession. Arrive with valid government-issued photo identification along with proof of funds, which could include a recent bank statement or a cashier's check request. If you intend to finance the purchase, discuss the final price independently of the loan documents. Once the negotiation is settled, present the precise sum of cash to the salesperson, who should sign a receipt confirming the payment. Immediately after the sale concludes, verify that the vehicle's title has been officially transferred into your name to prevent any subsequent legal complications regarding ownership.

\paragraph{Independent reference; quality passes; $p=5.22e-01$.} Before entering a dealership to purchase a vehicle, confirm that you possess a verified cash deposit well in advance of the visit. You must present government-issued ID and documentation of liquid assets, such as an updated bank record or a cashier's check, in order to proceed with the transaction. If your intention is to utilize debt coverage rather than paying cash directly, discuss the agreed-upon price with your loan documentation beforehand and then hand over the precise monetary value requested by the sales staff at the counter. Furthermore, require them to issue a receipt for the cash settlement and immediately confirm that the title deed has been legally transferred to your personal record to prevent any subsequent legal complications regarding the asset ownership.
